%% LyX 2.5.1 created this file.  For more info, see https://www.lyx.org/.
%% Do not edit unless you really know what you are doing.
\documentclass[brazil]{article}
\usepackage[T1]{fontenc}
\usepackage[utf8]{luainputenc}
\usepackage{geometry}
\geometry{verbose,tmargin=3cm,bmargin=3cm,lmargin=3cm,rmargin=2.5cm}
\usepackage{babel}
\begin{document}
\title{Introdução}

\maketitle
Redes neurais artificiais compõem uma classe de modelos computacionais que alcançaram uma alta taxa de sucesso em tarefas específicas envolvendo reconhecimento de padrões. Em aplicações práticas atuais, redes neurais artificiais são executadas usualmente como algoritmos clássicos em computadores convencionais \cite{Italiano}, mas há um crescente aumento de interesse em sua implementação em dispositivos quânticos. Propriedades intrínsecas da Mecânica Quântica como representação e armazenamento de grande quantidade de dados como vetores e matrizes, além da realização de operações lineares resultam em um aumento exponencial tanto na capacidade de armazenamento como no poder de processamento \cite{Do_Russo_1},\cite{Do_Russo_2},\cite{Do_Russo_3},\cite{Do_Russo_4}.

Entre os modelos de neurônios mais simples de redes neurais temos o perceptron. O perceptron foi proposto como um modelo probabilístico para o armazenamento e organização da informação no cérebro, descrito originalmente como uma aproximação frutífera na compreensão e organização dos sistemas cognitivos, permitia predizer as curvas de aprendizagem a partir de variáveis neurológicas e vice-versa \cite{rosenblatt1958perceptron}. Apesar de ser o modelo mais simples, ainda assim redes neurais construídas a partir de perceptrons, sob certos cenários, têm demonstrado benefícios quando comparado com técnicas de modelagem estatísticas mais tradicionais. Dentre estes benefícios, podemos destacar que, ao contrário de outras técnicas estatísticas, não faz suposições anteriores a distribuições de dados, podendo modelar funções altamente não lineares e pode ser treinado para generalizar com precisão quando defrontado com novos dados \cite{gardner1998artificial}.

Avançando ainda dentro das redes neurais, posteriormente também foram propostos modelos quânticos inspirados por modelos de redes neurais biológicas no qual os neurônios são sensíveis não somente a amplitude dos sinais de entrada, como é o caso do perceptron, mas também a fase. Desta maneira, uma evolução natural esperada dos modelos mais simples foi levar em conta a natureza complexa dos valores de entrada, e utilizar então a fase destes números para codificar a informação, além da amplitude até então utilizada \cite{Russo}, \cite{Esloveno}. Sutherland inclusive interpreta a fase como a informação em si, e a amplitude como o nível correspondente de confiança na mesma \cite{Do_Esloveno}.

Redes neurais foram responsáveis por avanços significativos e sem precedentes, gerando um enorme impacto no desenvolvimento da área. Porém, os métodos de kernel, apesar de terem ganho menos publicidade, também foram extensivamente estudados, tanto teoricamente quanto empiricamente, devido ao seu alto poder de modelagem de dados não-lineares. Uma das razões comumente atribuídas para isso, é a crença que seria difícil para os métodos de kernel acompanharem as redes neurais artificiais em problemas de larga escala, pois muitos algoritmos kernel escalam quadraticamente de acordo com o número de amostras de treinamento \cite{lu2014scale}. Porém este é um problema que pode ser superado com o emprego da computação quântica, uma vez que a exploração de espaços quânticos, que aumentam exponencialmente a capacidade de representação dos dados, é um dos elementos centrais desta mesma.

Dessa forma, aprendizado de máquina e computação quântica são duas tecnologias que possuem o potencial de alterar como a computação é feita e resolver problemas até então sem solução em tempo hábil. O uso da computação quântica em modelos de aprendizado de máquina se apresenta como uma possível aplicação em que há ganho de performance significativo se comparado com o melhor algoritmo possível em um computador clássico, isto é, o algoritmo quântico executado em um computador quântico é capaz de resolver determinados problemas em um tempo significativamente menor se comparado a alternativas clássicas \cite{havlivcek2019supervised}. Sendo assim, a computação quântica é uma área de intersecção entre a ciência da computação, física e a engenharia, que promete uma revolução na performance computacional devido ao seu paralelismo intrínseco, apesar de poder ter sido considerada uma utopia durante uma época, hoje em dia é uma realidade, uma vez que computadores quânticos reais podem ser acessados e programados através da internet por qualquer um \cite{acampora2019quantum}.

Atualmente o interesse por esta área não está mais restrito à academia. Grandes agências e empresas de reconhecimento internacional no setor tecnológico possuem projetos de conhecimento público na área. Como exemplo além da IBM, que disponibiliza dispositivos de hardware quântico para acesso via nuvem, podemos citar a NASA, a Amazon e o Google, sendo que este último desenvolve pesquisas exatamente na área de aprendizado de máquina em computadores quânticos \cite{moret2015can}. É fato que nos dias de hoje já temos um número considerável de algoritmos quânticos bem estabelecidos que apresentam um ganho de velocidade impressionante em comparação com os melhores métodos clássicos conhecidos. A computação quântica já é uma tecnologia emergente, e dentro desse cenário, pesquisas na área de aprendizado de máquina quântica parecem ser um caminho promissor para futuramente oferecer uma revolução nos métodos inteligentes de processamento de dados \cite{schuld2015introduction}.

Este trabalho encontra-se dividido da seguinte forma: Na próxima seção os objetivos são apresentados, na seção 3 é realizada uma revisão bibliográfica que cobre as teorias clássicas de aprendizado de máquina utilizadas para o desenvolvimento do resto do trabalho, computação e algoritmos quânticos, as adaptações necessárias nos modelos clássicos para serem executados em computadores quânticos e por fim uma rápida revisão sobre trabalhos relacionados. A seção 4 contém os resultados obtidos durante o desenvolvimento deste trabalho com as devidas discussões e o último capítulo trata-se então da conclusão.

\section{Objetivos}
\begin{itemize}
\item Objetivo acadêmico: Este trabalho tem como objetivo possibilitar a revisão e solidificação do conhecimento adquirido ao longo do curso sobre Mecânica Quântica, além de também servir como introdução à computação quântica e às teorias clássicas de aprendizado de máquina. 
\item Objetivo geral: O objetivo geral do presente trabalho é propor modelos eficientes de aprendizado de máquina para computadores quânticos e implementá-los tanto quanto possível, tendo em vista o atual acesso limitado à dispositivos quânticos. 
\item Objetivos específicos: 
\begin{itemize}
\item Implementar um modelo de neurônio binário baseado nos modelos de McCulloch-Pitts e Rosenblatt; 
\item Propor e implementar um modelo de neurônio contínuo baseado no modelo sigmoide; 
\item Propor e implementar um modelo baseado no método kernel de aprendizado de máquina. 
\end{itemize}
\end{itemize}
 \bibliographystyle{plain}
\bibliography{referencias}

\end{document}
